\section{Introduction}
\label{sec:intro}

GeoSection \citep{b8paper} introduced isoperimetrically-normalised ratio-optimal
bisection for congressional redistricting, demonstrating that minimum-edge-cut
recursive bisection with a $\sqrt{\min(i,k-i)}$ normalisation selects a
\emph{natural ratio} for each state that reflects genuine geographic structure
rather than computational artefact.
B.7 \citep{b7paper} established that this natural ratio is \emph{seed-stable}
within a single census year: running GeoSection with 1000 different random seeds
produces the same natural ratio in $>95\,\%$ of runs for most states.

This paper asks a deeper question: is the natural ratio \emph{census-stable}?

\paragraph{The census-stability question.}
Every ten years, the Census Bureau publishes new population data, tract boundary
updates, and a new congressional apportionment.
Redistricting commissions and courts must draw new district maps from scratch.
The question for an algorithmic redistricting system is: if we run GeoSection on
2000 Census data, then on 2010 data, then on 2020 data, do we get the same
answer?

If yes, the result is remarkable.
The 2000 census measured one population distribution.
The 2010 census measured a different distribution---ten years of births, deaths,
migration, and urban development later.
The 2020 census measured yet another.
Despite these three independent snapshots of American demographic reality, the
algorithm converges to the same geographic partition.
That partition is not a product of any particular census cycle.
It is the \emph{geographic structure of the state}.

\paragraph{Why this matters legally.}
The most powerful defence against a partisan gerrymandering claim is
demonstrating that the challenged map was not drawn to achieve a particular
partisan outcome.
B.7's seed stability provides one dimension of this defence: the map is the
same regardless of which random seed initialised the algorithm.
Census stability provides a second, stronger dimension.

Consider a state where GeoSection has produced the same 5\,D/9\,R seat
distribution in 2000, 2010, and 2020.
A challenger who claims the map was drawn to achieve Republican advantage faces
a severe epistemic problem: the algorithm produced the same answer under three
different population datasets spanning twenty years.
The $-13$\,pp proportionality gap is not a product of 2020's specific partisan
distribution.
It was the result in 2000, when Republicans had a different geographic
distribution.
It was the result in 2010, after Obama-era demographic shifts.
It was the result in 2020.
The geographic sorting of Democratic voters into concentrated urban areas is
the cause of the gap---and that sorting is stable across census cycles.

In state-court litigation under post-\emph{Rucho} doctrine---where courts must
find that a deviation from the natural map cannot be explained by
geography---a census-stable GeoSection map carries the strongest possible
geographic determinism claim:
\emph{this partition has existed before living memory of the current partisan
alignment}.

\paragraph{The Census Stability Score.}
We formalise this intuition as the \textbf{Census Stability Score} (CSS), a
composite metric computed for each state by comparing GeoSection outputs across
census years:
\[
  \mathrm{CSS}(\text{state}) = 0.5 \cdot s_{\text{seat}}
                              + 0.3 \cdot s_{\text{ratio}}
                              + 0.2 \cdot s_{\text{gap}},
\]
where $s_{\text{seat}}$ measures partisan-outcome stability,
$s_{\text{ratio}}$ measures natural-ratio stability, and $s_{\text{gap}}$
measures proportionality-gap-magnitude stability.
A CSS of 1.0 means identical seat counts, identical natural ratios, and
identical proportionality gaps across all census years.
CSS $\geq 0.90$ is the threshold for ``highly stable'': the state's geographic
structure is durably expressed by GeoSection across decades.

\paragraph{Two sources of change.}
A map may change between census years for two reasons.
\emph{Type I}: population shifts within a fixed geographic topology.
People move between tracts, altering population balance targets, but the
underlying tract graph is unchanged.
This is a smooth perturbation of the METIS input weights.
\emph{Type II}: tract boundary redesign.
The Census Bureau redraws some tract boundaries between decennial censuses,
splitting or merging tracts.
This changes the graph topology and can propagate to downstream district
boundaries.
StabilitySection separates these two effects, measuring each independently.

\paragraph{Relationship to B.7 and B.8.}
B.7 measures stability along the \emph{seed dimension}: fixed census year,
varying initialisation.
B.15 measures stability along the \emph{census dimension}: fixed algorithm,
varying input data.
Together they provide a 2D stability picture:
$(s_\text{seed}, s_\text{census})$.
States with high stability in both dimensions have the strongest geographic
determinism claim.
States with low stability in both dimensions require the most careful
procedural justification for any particular map they produce.

\paragraph{Our contributions.}
\begin{enumerate}
\item \textbf{The Census Stability Score (CSS)} (\S\ref{sec:methodology}):
  a formally defined composite metric for cross-census algorithmic stability,
  with sub-components measuring seat, ratio, and gap stability separately.

\item \textbf{Lorenz drift as a predictive proxy} (\S\ref{sec:lorenz}):
  the shift in the optimal ratio $p^*$ on the population-geography Lorenz
  curve between census years predicts CSS; states where $p^*$ drifts
  significantly are structurally volatile.

\item \textbf{Two-source decomposition} (\S\ref{sec:decomposition}):
  population-only perturbation stability vs.\ full cross-year stability,
  isolating the contribution of tract-boundary redesign.

\item \textbf{50-state 2020 GeoSection baseline} (\S\ref{sec:evaluation}):
  complete natural ratio and proportionality results for 44 multi-district
  states, providing the 2020 anchor for the cross-census comparison.

\item \textbf{Legal and policy implications} (\S\ref{sec:legal}):
  the CSS as a litigation tool and proposed ``Districting Integrity'' statute
  design.
\end{enumerate}

\paragraph{Data status.}
The 2020 GeoSection sweep (44 states, 50 seeds per ratio) is complete.
The 2000 census sweep has completed 47 of 50 states.
The 2010 census sweep is in progress.
\textbf{[TBD: 2010 sweep pending]}: the three-census CSS computation awaits
2010 completion.
The two-census comparison (2000 vs.\ 2020) is reported in full.

\paragraph{Paper organisation.}
Section~\ref{sec:related} reviews related work.
Section~\ref{sec:methodology} defines StabilitySection and the CSS formally.
Section~\ref{sec:evaluation} presents the 2020 baseline and available
cross-census comparison.
Section~\ref{sec:discussion} discusses legal implications and predictive
modelling.
Section~\ref{sec:conclusion} concludes.
